% !TeX program = pdflatex
\documentclass[aspectratio=169,xcolor={dvipsnames,table}]{beamer}

\input{../../_en_preambulo_beamer}
% Comandos y macros estadísticas compartidas para las presentaciones Beamer en inglés (Metropolis)

% Conjuntos numéricos básicos
\newcommand{\R}{\mathbb{R}}
\newcommand{\N}{\mathbb{N}}
\newcommand{\Q}{\mathbb{Q}}
\newcommand{\Z}{\mathbb{Z}}
\newcommand{\set}[1]{\left\{ #1 \right\}}
\newcommand{\abs}[1]{\left|#1\right|}
\newcommand{\norm}[1]{\left\| #1 \right\|}

% Comandos estadísticos y probabilísticos del curso
\newcommand{\Var}{\operatorname{Var}}
\newcommand{\cov}{\operatorname{Cov}}
\newcommand{\corr}{\rho}
\newcommand{\comb}[2]{\begin{pmatrix} #1 \\ #2 \end{pmatrix}}
\newcommand{\s}{\sigma}
\newcommand{\particion}{\mathfrak{P}}
\newcommand{\card}[1]{\#\left( #1 \right)}
\newcommand{\seq}[3]{\set{#1}_{#2}^{#3}}
\newcommand{\E}{\mathbb{E}}
\newcommand{\Prob}{\mathbb{P}}

% Entornos pedagógicos y cajas en inglés académico natural (soportando tanto nombres en inglés como en español)
\newenvironment{discussion}{%
  \begin{alertblock}{Discussion Question}%
}{%
  \end{alertblock}%
}
\newenvironment{discusion}{%
  \begin{alertblock}{Discussion Question}%
}{%
  \end{alertblock}%
}

\newenvironment{reflection}{%
  \begin{alertblock}{Classroom Reflection}%
}{%
  \end{alertblock}%
}
\newenvironment{reflexion}{%
  \begin{alertblock}{Classroom Reflection}%
}{%
  \end{alertblock}%
}

\newenvironment{paradox}{%
  \begin{alertblock}{Exploratory Paradox}%
}{%
  \end{alertblock}%
}
\newenvironment{paradoja}{%
  \begin{alertblock}{Exploratory Paradox}%
}{%
  \end{alertblock}%
}

\newenvironment{motivation}{%
  \begin{exampleblock}{Motivation: Data Science in Practice}%
}{%
  \end{exampleblock}%
}
\newenvironment{motivacion}{%
  \begin{exampleblock}{Motivation: Data Science in Practice}%
}{%
  \end{exampleblock}%
}

\newenvironment{quickchallenge}{%
  \begin{block}{Quick Team Challenge}%
}{%
  \end{block}%
}
\newenvironment{retoclase}{%
  \begin{block}{Quick Team Challenge}%
}{%
  \end{block}%
}


\title{Central Limit Theorem: Asymptotic Convergence}
\subtitle{Section 05.02 --- MGF Proof, Convergence in Distribution, and Berry-Esseen}
\author[J. Castillo Colmenares]{Juliho Castillo Colmenares}
\institute[Tec de Monterrey]{Tecnologico de Monterrey}
\date{\vspace{-1.5cm}}

\begin{document}

% SLIDE 01: Title Page
\begin{frame}[plain]
	\titlepage
\end{frame}

% SLIDE 02: Roadmap
\begin{frame}{Roadmap --- Chapter 05: Sampling Distributions}
	\vspace{-0.15cm}
	\begin{block}{Curricular Structure of Unit 4}
		\footnotesize
		\begin{enumerate}\itemsep=0.03cm
			\item \textcolor{gray}{Section 05.01: Simple Random Sampling, Mean, and Unbiased Sample Variance}
			\item \textbf{\textcolor{TecRojo}{Section 05.02: Asymptotic Central Limit Theorem}} $\leftarrow$ \emph{Current focus}
			\item \textcolor{gray}{Section 05.03: Chi-Squared Distribution and Sample Variance}
			\item \textcolor{gray}{Section 05.04: Student's $t$-Distribution and Small Samples}
			\item \textcolor{gray}{Section 05.05: Fisher-Snedecor $F$-Distribution}
		\end{enumerate}
	\end{block}
	\vspace{-0.2cm}
	\begin{alertblock}{Learning Objective}
		\scriptsize
		Formalize the Central Limit Theorem through convergence in distribution, prove it rigorously via the moment-generating function, and quantify its rate of convergence with the Berry-Esseen theorem, beyond the ``$n\ge30$'' heuristic.
	\end{alertblock}
\end{frame}

% SLIDE 03: Motivation
\begin{frame}{How Fast Does the CLT Converge?}
	\footnotesize
	\begin{motivacion}
		In Section 02.06 we saw the CLT introductorily: the sample mean approaches a Normal when $n$ is ``large enough.'' But what exactly does ``large enough'' mean? Why does it work regardless of the original distribution? Today we formalize both questions.
	\end{motivacion}
	\vspace{0.15cm}
	\pause
	\begin{itemize}\itemsep=0.1cm
		\item \textbf{Convergence in distribution:} the precise type of limit that applies to the CLT. \pause
		\item \textbf{MGF proof:} why the Normal emerges as a universal limit. \pause
		\item \textbf{Berry-Esseen:} an explicit bound on the approximation error for finite $n$.
	\end{itemize}
\end{frame}

% SLIDE 03B: Convergence in Distribution and CLT
\begin{frame}{Convergence in Distribution and the Central Limit Theorem}
	\vspace{-0.15cm}
	\begin{columns}[T]
		\begin{column}{0.48\textwidth}
			\begin{block}{Convergence in Distribution}
				\footnotesize
				$Y_n \xrightarrow{d} Y$ if $F_{Y_n}(y)\to F_Y(y)$ at every continuity point of $F_Y$.
			\end{block}
		\end{column}
		\begin{column}{0.48\textwidth}
			\begin{alertblock}{CLT (Lindeberg-Levy)}
				\footnotesize
				For iid $X_i$ with mean $\mu$, variance $\sigma^2<\infty$: $Z_n=(\bar X_n-\mu)/(\sigma/\sqrt n) \xrightarrow{d} N(0,1)$ as $n\to\infty$, regardless of the original distribution.
			\end{alertblock}
		\end{column}
	\end{columns}
\end{frame}

% SLIDE 04: Proof via MGF
\begin{frame}{Proof of the CLT via the Moment-Generating Function}
	\vspace{-0.15cm}
	\begin{columns}[T]
		\begin{column}{0.48\textwidth}
			\begin{block}{Setup}
				\footnotesize
				$M_{Z_n}(t) = e^{-\mu t\sqrt n/\sigma}[M_{X_1}(t/(\sigma\sqrt n))]^n$. Expand $\log M_{X_1}(s) = \mu s + \sigma^2 s^2/2 + O(s^3)$.
				\end{block}
		\end{column}
		\begin{column}{0.48\textwidth}
			\begin{alertblock}{Taking the Limit}
				\footnotesize
				$\log M_{Z_n}(t) = t^2/2 + O(n^{-1/2}) \to t^2/2$. By Levy's continuity theorem, $Z_n \xrightarrow{d} N(0,1)$.
			\end{alertblock}
		\end{column}
	\end{columns}
\end{frame}

% SLIDE 05: Berry-Esseen
\begin{frame}{Berry-Esseen: Rate of Convergence}
	\vspace{-0.15cm}
	\begin{columns}[T]
		\begin{column}{0.48\textwidth}
			\begin{block}{Berry-Esseen Bound}
				\footnotesize
				With $\rho=E|X_i-\mu|^3<\infty$ and universal $C<0.5$: $\sup_z|F_{Z_n}(z)-\Phi(z)| \le C\rho/(\sigma^3\sqrt n)$.
			\end{block}
		\end{column}
		\begin{column}{0.48\textwidth}
			\begin{alertblock}{The ``$n\ge30$'' Rule in Context}
				\footnotesize
				Error decays as $O(1/\sqrt n)$: quadrupling $n$ halves the max error. Heavily skewed populations (large $\rho$) need larger $n$.
			\end{alertblock}
		\end{column}
	\end{columns}
\end{frame}

% SLIDE 06: Applications
\begin{frame}{Applications: Sums and Proportions}
	\vspace{-0.1cm}
	\begin{itemize}\itemsep=0.1cm
		\item \textbf{Sums:} if $T=\sum X_i$, then $T \approx N(n\mu, n\sigma^2)$ (insurance, inventory control). \pause
		\item \textbf{Proportions:} $\hat p = \bar X$ for $X_i \sim \text{Bernoulli}(p)$, with $\hat p \approx N(p, p(1-p)/n)$ (surveys, A/B testing). \pause
		\item \textbf{Binomial-Normal approximation:} a special case of sums of Bernoullis, with Yates continuity correction. \pause
		\item \textbf{Bootstrap resampling:} the CLT justifies why bootstrap distributions approach normality for large $n$.
	\end{itemize}
\end{frame}

% SLIDE 07: Python Lab Block 1
\begin{frame}[fragile]{Python Lab: Convergence from a Skewed Population}
	\vspace{-0.05cm}
	\scriptsize
	Kolmogorov-Smirnov test of $Z_n$ against $N(0,1)$ for an Exponential population (\texttt{05.02\_central\_limit\_theorem.py}):
	{\lstset{basicstyle=\fontsize{4pt}{4.8pt}\selectfont\ttfamily,aboveskip=0.1em,belowskip=0.1em}
	\lstinputlisting[firstline=17, lastline=35]{../../code/05_distribuciones_muestreo/05.02_central_limit_theorem.py}}
\end{frame}

% SLIDE 08: Python Lab Block 2
\begin{frame}[fragile]{Python Lab: Berry-Esseen Rate of Convergence}
	\vspace{-0.05cm}
	\scriptsize
	Empirical estimate of the maximum gap $\sup_z|F_{Z_n}(z)-\Phi(z)|$ and verification of the $O(1/\sqrt n)$ rate:
	{\lstset{basicstyle=\fontsize{4pt}{4.8pt}\selectfont\ttfamily,aboveskip=0.1em,belowskip=0.1em}
	\lstinputlisting[firstline=38, lastline=64]{../../code/05_distribuciones_muestreo/05.02_central_limit_theorem.py}}
\end{frame}

% SLIDE 09: Python Lab Block 3
\begin{frame}[fragile]{Python Lab: CLT for Sums and Proportions}
	\vspace{-0.05cm}
	\scriptsize
	Applying the CLT to a total insurance claim amount and to a sample proportion, with Monte Carlo verification:
	{\lstset{basicstyle=\fontsize{4pt}{4.8pt}\selectfont\ttfamily,aboveskip=0.1em,belowskip=0.1em}
	\lstinputlisting[firstline=67, lastline=96]{../../code/05_distribuciones_muestreo/05.02_central_limit_theorem.py}}
\end{frame}

% SLIDE 10: Python Lab Terminal Output
\begin{frame}[fragile]{Python Lab: Terminal Output}
	\vspace{-0.1cm}
	\begin{block}{}
		\fontsize{4pt}{5pt}\selectfont
		\begin{verbatim}
=== Block 1: CLT Convergence from a Skewed Population ===
Exp(lambda=0.25): n=5 D=0.0612 | n=30 D=0.0253 | n=100 D=0.0141 (D shrinks with n)

=== Block 2: Berry-Esseen Rate of Convergence ===
n=10: max gap=0.0418 | n=40: 0.0220 | n=160: 0.0120
Ratio n=10->40: 0.526 (theo 0.500) | n=40->160: 0.545 (theo 0.500)

=== Block 3: CLT for Sums and Proportions ===
Insurance: E(T)=80000, SD(T)=3000, Z=1.667, P(T>85000)=0.0478
Proportion: SD(phat)=0.0324, Z=1.543, P(phat>0.35)=0.0614 (MC: 0.0541)
		\end{verbatim}
	\end{block}
\end{frame}

% SLIDE 11: Exercise Level 1 (Statement)
\begin{frame}{In-Class Exercise --- Level 1: Fundamental (1/2)}
	\vspace{-0.1cm}
	\begin{block}{Problem 5.2.1 --- Direct CLT Application (Statement)}
		A population has mean $\mu=200$ and standard deviation $\sigma=50$. A simple random sample of $n=100$ is drawn. Use the CLT to approximate $P(180 < \bar{X} < 210)$.
	\end{block}
	\vspace{0.15cm}
	\pause
	\begin{alertblock}{Model Setup}
		\begin{itemize}\itemsep=0.04cm
			\item $\text{SD}(\bar{X}) = \sigma/\sqrt{n} = 50/10 = 5$.
			\item Standardize both endpoints of the interval.
		\end{itemize}
	\end{alertblock}
\end{frame}

% SLIDE 12: Exercise Level 1 (Solution)
\begin{frame}{In-Class Exercise --- Level 1: Fundamental (2/2)}
	\vspace{-0.05cm}
	\scriptsize
	Standardize both endpoints: \pause
	\begin{block}{Calculation}
		\begin{align*}
			P(180 < \bar{X} < 210) \approx P\left(\dfrac{180-200}{5} < Z < \dfrac{210-200}{5}\right) = P(-4<Z<2) = \Phi(2)-\Phi(-4) \approx 0.9772.
		\end{align*}
	\end{block}
	\vspace{0.05cm}
	\pause
	\begin{alertblock}{Interpretation}
		\scriptsize
		Since $\Phi(-4)\approx 0$, virtually all the relevant probability comes from the upper tail: there is about a $97.7\%$ probability that the sample mean falls in the given range.
	\end{alertblock}
\end{frame}

% SLIDE 13: Exercise Level 2 (Statement)
\begin{frame}{In-Class Exercise --- Level 2: Operational (1/2)}
	\vspace{-0.1cm}
	\begin{block}{Problem 5.2.4 --- CLT for a Sum (Statement)}
		An insurer receives $n=100$ independent claims, each with mean $\mu=\$800$ and standard deviation $\sigma=\$300$. Use the CLT to approximate the probability that the total amount $T=\sum X_i$ exceeds $\$85{,}000$.
	\end{block}
	\vspace{0.15cm}
	\pause
	\begin{alertblock}{Strategy}
		\scriptsize
		For a sum (not an average), use $E(T)=n\mu$ and $\Var(T)=n\sigma^2$ (do not divide by $n$).
	\end{alertblock}
\end{frame}

% SLIDE 14: Exercise Level 2 (Solution)
\begin{frame}{In-Class Exercise --- Level 2: Operational (2/2)}
	\vspace{-0.05cm}
	\scriptsize
	Compute the mean and standard deviation of the sum: \pause
	\begin{block}{Calculation}
		\begin{align*}
			E(T) = 100\cdot 800 = 80{,}000, \qquad \text{SD}(T)=\sqrt{100}\cdot 300 = 3000.
		\end{align*}
		\begin{align*}
			P(T>85{,}000) \approx P\left(Z>\dfrac{85{,}000-80{,}000}{3000}\right) = P(Z>1.667) \approx 0.0478.
		\end{align*}
	\end{block}
	\vspace{0.05cm}
	\pause
	\begin{alertblock}{Interpretation}
		\scriptsize
		There is roughly a $4.78\%$ probability that total claims exceed the $\$85{,}000$ budget; the insurer can use this to size capital reserves.
	\end{alertblock}
\end{frame}

% SLIDE 15: Exercise Level 3 (Statement)
\begin{frame}{In-Class Exercise --- Level 3: Analytical (1/2)}
	\vspace{-0.1cm}
	\begin{block}{Problem 5.2.7 --- CLT for the Exponential via MGF (Statement)}
		Verify the MGF proof of the CLT specifically for $X_i \sim \text{Exp}(\lambda)$ (with $\mu=1/\lambda$, $\sigma=1/\lambda$), using $M_{X_1}(t)=(1-t/\lambda)^{-1}$.
	\end{block}
	\vspace{0.1cm}
	\pause
	\begin{alertblock}{Strategy}
		\scriptsize
		Substitute $\mu=\sigma=1/\lambda$ into the general formula for $M_{Z_n}(t)$; note that $\lambda\sigma=1$ simplifies the expression considerably.
	\end{alertblock}
\end{frame}

% SLIDE 16: Exercise Level 3 (Solution)
\begin{frame}{In-Class Exercise --- Level 3: Analytical (2/2)}
	\vspace{-0.05cm}
	\scriptsize
	Substitute and simplify using $\lambda\sigma=1$: \pause
	\begin{block}{Proof}
		\scriptsize
		\begin{align*}
			M_{Z_n}(t) = e^{-t\sqrt{n}}\left(1-\dfrac{t}{\sqrt{n}}\right)^{-n}.
		\end{align*}
		With $-\log(1-t/\sqrt n) = t/\sqrt n + t^2/(2n)+O(n^{-3/2})$:
		\begin{align*}
			\log M_{Z_n}(t) = -t\sqrt{n} + n\left[\dfrac{t}{\sqrt{n}}+\dfrac{t^2}{2n}+O(n^{-3/2})\right] = \dfrac{t^2}{2}+O(n^{-1/2}) \to \dfrac{t^2}{2}. \quad \qedhere
		\end{align*}
	\end{block}
	\vspace{0.05cm}
	\pause
	\begin{alertblock}{Interpretation}
		\scriptsize
		Even though the Exponential is strongly skewed, the CLT still converges: the general proof does not require symmetry, only finite mean and variance.
	\end{alertblock}
\end{frame}

% SLIDE 17: Exercise Level 4 (Statement)
\begin{frame}{In-Class Exercise --- Level 4: Challenging (1/2)}
	\vspace{-0.1cm}
	\begin{block}{Problem 5.2.10 --- CLT with the Finite Population Correction (Statement)}
		A finite population of $N=500$ accounts has $\sigma=\$80$. A sample of $n=50$ accounts is drawn \emph{without replacement}. Combine the CLT with the FPC (Section 05.01) to approximate $P(\bar{X} > \mu + 15)$.
	\end{block}
	\vspace{0.15cm}
	\pause
	\begin{alertblock}{Strategy}
		\scriptsize
		First compute $\Var(\bar X)$ with the FPC, then apply the CLT using that adjusted variance instead of $\sigma^2/n$.
	\end{alertblock}
\end{frame}

% SLIDE 18: Exercise Level 4 (Solution)
\begin{frame}{In-Class Exercise --- Level 4: Challenging (2/2)}
	\vspace{-0.05cm}
	\scriptsize
	Combine the FPC (05.01) with today's CLT: \pause
	\begin{block}{Calculation}
		\scriptsize
		\begin{align*}
			\Var(\bar{X}) = \dfrac{80^2}{50}\cdot\dfrac{450}{499} = 128\cdot 0.9018 \approx 115.43, \qquad \text{SD}(\bar{X})\approx 10.744.
		\end{align*}
		\begin{align*}
			P(\bar{X}>\mu+15) \approx P\left(Z>\dfrac{15}{10.744}\right) = P(Z>1.396) \approx 0.0814.
		\end{align*}
	\end{block}
	\vspace{0.05cm}
	\pause
	\begin{alertblock}{Interpretation}
		\scriptsize
		This problem connects two sections: the FPC adjusts the variance for sampling without replacement from a finite population, and the CLT lets us use the Normal distribution to compute the final probability, even without knowing the exact distribution of the accounts.
	\end{alertblock}
\end{frame}

% SLIDE 19: Closing and Chapter Perspective
\begin{frame}{Closing Section and Chapter Perspective}
	\vspace{-0.2cm}
	\begin{block}{Synthesis: The Central Limit Theorem, Formalized}
		\scriptsize
		The Central Limit Theorem is, in a precise sense, the most important result in classical statistics: it guarantees that the sample mean converges to a Normal distribution regardless of the shape of the original population. Today we formalized this result with a rigorous MGF-based proof and quantified its rate of convergence via the Berry-Esseen theorem.
	\end{block}
	\vspace{0.05cm}
	\pause
	\begin{alertblock}{Key Lessons and Next Step}
		\begin{itemize}\itemsep=0.05cm
			\item \textbf{CLT:} $Z_n = (\bar{X}_n-\mu)/(\sigma/\sqrt{n}) \xrightarrow{d} N(0,1)$, regardless of the original distribution.
			\item \textbf{MGF proof:} $\log M_{Z_n}(t) \to t^2/2$ for any population with finite moments.
			\item \textbf{Berry-Esseen:} the approximation error decays as $O(1/\sqrt{n})$, quantifying the ``$n\ge30$'' heuristic.
			\item \textbf{Next Section:} \textcolor{TecRojo}{Chi-Squared Distribution and Sample Variance}.
		\end{itemize}
	\end{alertblock}
\end{frame}

\end{document}
